\begin{abstract}
Chunking ảnh hưởng trực tiếp đến khả năng truy xuất bằng chứng và chất lượng câu trả lời trong hệ thống \textit{Retrieval-Augmented Generation} (RAG). Nghiên cứu này so sánh ba chiến lược \textit{fixed-size}, \textit{structure-aware} và \textit{prompt-based} trên 384 tài liệu kỹ thuật Angular ở định dạng Markdown. Nhóm đồng thời xây dựng một bộ benchmark gồm 140 câu hỏi gắn với bằng chứng nguồn và đánh giá các chiến lược trong cùng một pipeline được kiểm soát. Kết quả cho thấy \textit{fixed-size} phù hợp hơn khi ưu tiên thứ hạng truy xuất và chất lượng câu trả lời, trong khi \textit{structure-aware} có lợi thế về khả năng thu hồi đầy đủ bằng chứng. \textit{Prompt-based} chưa thể hiện ưu thế nhất quán. Đáng chú ý, khả năng thu hồi bằng chứng tốt hơn không nhất thiết chuyển thành chất lượng câu trả lời cao hơn. Vì vậy, chiến lược chunking cần được lựa chọn theo mục tiêu của hệ thống RAG thay vì dựa trên một chỉ số duy nhất.
\end{abstract}